\begin{table}[H]
\centering
\caption{Inheritance Systems and Birth-Order Productivity After the Sixteenth Century}
\label{tab:inheritance_post15}

\begin{tablebody}
\begin{tabular}{l c}
\toprule
 & (1) \\
 & \textit{Asinh difference in titles} \\
\midrule

Intercept
& 0.998*** \\
& (0.305) \\

Non-egalitarian inheritance
& -0.362 \\
& (0.631) \\

Inheritance unknown
& -1.716 \\
& (1.258) \\

Post-16th century
& -1.134*** \\
& (0.314) \\

Age gap
& 0.019* \\
& (0.010) \\

Father in academia
& -0.178 \\
& (0.149) \\

Different institution
& 0.472*** \\
& (0.180) \\

Non-egalitarian $\times$ Post-16th
& 0.533 \\
& (0.655) \\

Unknown inheritance $\times$ Post-16th
& 1.897 \\
& (1.276) \\

\midrule
Clustered SE (family) & Yes \\
\midrule
Observations & 645 \\
R$^{2}$ & 0.034 \\
Adj. R$^{2}$ & 0.022 \\
\bottomrule
\end{tabular}
\end{tablebody}

\vspace{4pt}
\begin{minipage}{\textwidth}
\footnotesize
\textit{Notes:} The dependent variable is the inverse hyperbolic sine (asinh) of the difference in titles between older and younger brothers. The specification interacts inheritance systems with an indicator for observations from the sixteenth century onward. The split is motivated by the limited variation in inheritance-system classifications observed before the sixteenth century. Standard errors are clustered at the family level. \textit{Significance levels:} *** $p<0.01$, ** $p<0.05$, * $p<0.10$.
\end{minipage}

\end{table}